% Failure settings (generated from results/stress_runs.csv; see results/summary.md, "Stress study").
In the main sweeps robots almost never fail, but partly by construction: the scenario
generator removes disruptions that would make a task impossible (Section 2). To find
where the agents really fail, we ran a \textbf{stress study} (760 runs, 10 seeds, all
strategies) in which this filter can be switched off and the environment pushed further.
\begin{itemize}
  \item \textbf{unsafe}: filter off, half of the blockages and breakdowns permanent, $\rho$ = 0--20\%, 20 robots;
  \item \textbf{unsafe-single}: one permanent blockage or breakdown with the filter off, 20 and 40 robots;
  \item \textbf{dense}: filter on, $\rho$ = 20--30\%, 20 robots;
  \item \textbf{crowd}: a small $12\times27$ map (204 free cells, 46 docks) with 20--44 robots;
  \item \textbf{radius}: communication radius $\R$ = 1, 2, 3, 6 with 40 robots.
\end{itemize}
A run \emph{fails} if some task is still undone when it ends. A run ends when every robot
is done, or when no robot has moved and no goal has been reached for 150 steps. Each
unfinished robot is classified as \emph{goal cut off} (a remaining goal is unreachable
past permanent obstacles and dead robots, so the instance has become unsolvable) or
\emph{deadlocked} (every goal is reachable, but the robot is stuck). No run had a collision.

\begin{figure}[H]
  \centering
  \includegraphics[width=\linewidth]{stress_completion.png}
  \caption{Share of tasks completed in the stress settings.}
  \label{fig:stress}
\end{figure}

\begin{table}[H]
\centering\footnotesize\setlength{\tabcolsep}{4pt}
\caption{Settings where agents fail. Runs failed out of 10 and the share of tasks
completed; stuck robots per run for \texttt{local} (cut off / deadlocked).}
\label{tab:failures}
\begin{tabular}{llrrrrr}
\toprule
setting & value & \multicolumn{2}{c}{\texttt{local} (v2)} & stuck & \texttt{local-flat} & \texttt{full} \\
 & & failed & completed & cut / dead & completed & completed \\
\midrule
unsafe & $\rho=0\%$  & 3/10 & 96.9\% & 0.5 / 0.2 & 96.9\% & 96.9\% \\
       & $\rho=5\%$  & 7/10 & 91.0\% & 1.5 / 0.3 & 91.0\% & 90.8\% \\
       & $\rho=10\%$ & 10/10 & 80.6\% & 2.3 / 3.1 & 79.2\% & 71.3\% \\
       & $\rho=20\%$ & 10/10 & 26.8\% & 9.9 / 8.8 & 27.4\% & 24.8\% \\
\midrule
unsafe-single & 1 perm.\ breakdown, 20 robots & 2/10 & 97.5\% & 0.4 / 0.2 & 97.5\% & 97.5\% \\
              & 1 perm.\ breakdown, 40 robots & 4/10 & 96.7\% & 0.4 / 1.0 & 96.7\% & 96.7\% \\
              & 1 perm.\ blockage, 20 or 40 robots & 1/10 & 99.7\% & 0.1 / 0.0 & 99.7\% & 99.7\% \\
\midrule
dense & $\rho=25\%$ & 0/10 & 100\% & -- & 96.3\% & 100\% \\
      & $\rho=30\%$ & 1/10 & 92.9\% & 0.0 / 2.0 & 92.7\% & 86.1\% \\
\midrule
crowd & 30 robots & 1/10 & 92.1\% & 0.0 / 2.9 & 92.1\% & 100\% \\
      & 40 robots & 3/10 & 73.9\% & 0.0 / 11.9 & 90.1\% & 94.9\% \\
      & 44 robots & 0/10 & 100\% & -- & 93.4\% & 100\% \\
\midrule
radius & $\R$ = 1, 2, 3, 6 & 0/10 & 100\% & -- & 100\% & 100\% \\
\bottomrule
\end{tabular}
\end{table}

\begin{figure}[H]
  \centering
  \includegraphics[width=0.95\linewidth]{stress_causes.png}
  \caption{Why robots fail to finish under \texttt{local}: goals cut off (red) or deadlocked
  although every goal is reachable (orange).}
  \label{fig:causes}
\end{figure}

\subsection*{F1. Disruptions that make tasks impossible}
With the filter off, a permanent blockage or a robot that breaks down for good can cut
off a pickup cell, a station or a whole aisle segment (Figure~\ref{fig:unreachable}). No
repair can complete those tasks, so all strategies fail alike. Completion falls from
96.9\% at $\rho=0\%$ to 26.8\% at $\rho=20\%$, where the map breaks into pieces and about
10 robots per run have a goal cut off. Even with no blockages at all ($\rho=0\%$), 3 of 10
runs fail, because a permanently broken robot in a one-cell aisle is itself a wall. A
\emph{single} permanent breakdown is enough to make 2 of 10 runs fail at 20 robots and
4 of 10 at 40.

\begin{figure}[H]
  \centering
  \includegraphics[width=0.85\linewidth]{failure_unreachable.png}
  \caption{Failure F1/F2 (\texttt{results/demo\_failure/unreachable.gif}). Permanent
  blockages (black) cut off the goals of robots 3, 13 and 18. They hold forever and
  5 of 62 tasks are never done.}
  \label{fig:unreachable}
\end{figure}

\subsection*{F2. Knock-on deadlock behind a stuck robot}
A robot whose goal is cut off does not give up its task. It holds its cell indefinitely,
and in one-cell aisles that cell can be the only way to another robot's goal. At
$\rho=10\%$, 2.3 robots per run have a goal cut off, but another 3.1 are deadlocked even
though all of their goals are reachable. At 20\% it is 9.9 and 8.8. More than half of the
failures in these settings are therefore caused by the agents' reaction, not by the
instance. Local repair suffers less from this than full replanning (80.6\% vs 71.3\%
completed at 10\%, 3.1 vs 6.6 deadlocked robots), because a global replan rebuilds every
plan around the stuck robots and can wall more robots in.

\subsection*{F3. Hold cascade in crowded maps (a weakness of our method)}
On the small map with 40 robots, 3 of 10 \texttt{local} runs fail, and in each of them
nearly the whole fleet stops. In seed 3, a robot breaks down \emph{temporarily} in an aisle
at $t=14$ and directly blocks 6 robots. They cannot find a detour, so they hold (Tier 4).
Every robot that planned through a held cell is then re-queued and also holds. Within one
step all 40 robots hold, each seeing the others as permanent walls, and only 1 of 122
tasks is ever delivered. Negotiation cannot untangle this: coalitions have at most 3--4
members, and Tier 3b at most 8, while the knot involves 40. \texttt{full} replanning plans
everyone together and solves the same instance (122/122). This is the setting in which
the local method is clearly worse than the global baseline. \texttt{local-flat} failed in
1 of 10 runs here and \texttt{solo} in 4. The effect depends strongly on the instance: at 44
robots no \texttt{local} run failed, and the 95\% confidence band in
Figure~\ref{fig:stress} is wide.

\subsection*{F4. Very high obstacle density}
With the filter on (every task stays solvable), failures start at about 25--30\% density:
1 of 10 runs fails for each local strategy at 30\% (2 deadlocked robots), and 2 of 10 for
full replanning (86.1\% completed). Up to 20\% density no strategy failed.

\subsection*{F5. Gridlock without negotiation}
Without negotiation (\texttt{solo}), robots gridlock much more often. On the main map with
50 robots, a single permanent breakdown at $t=40$ leaves all 49 remaining robots holding
from then on, and only 26 of 150 tasks are done (Figure~\ref{fig:gridlock}). On the small map
with 44 robots, 5 of 10 \texttt{solo} runs fail (59.7\% completed). \texttt{local} solves
both: 150/150 on the first instance and 10 of 10 runs on the second. Negotiation is what
prevents these gridlocks.

\begin{figure}[H]
  \centering
  \includegraphics[width=0.85\linewidth]{failure_gridlock_solo.png}
  \caption{Failure F5 (\texttt{results/demo\_failure/gridlock\_solo.gif}): with
  \texttt{solo}, all 49 live robots (red) gridlock after one permanent breakdown. The same
  instance with \texttt{local} completes all 150 tasks.}
  \label{fig:gridlock}
\end{figure}

\subsection*{F6. Emergency deadlines}
Emergency tasks are always eventually delivered, but rarely by their deadline (ideal time
+ 3 steps) once the fleet is dense: 25--35\% on time at 30--50 robots on the main map,
and 5\% on the crowded map with 40 robots and at 30\% density. Full replanning reaches
70--100\% because it can move every robot out of the emergency robot's way.

\subsection*{Where the agents did not fail}
\begin{itemize}
  \item \textbf{Communication radius.} Shrinking $\R$ from 6 to 1 caused no failures at
        40 robots and raised SoC by only 1\% (5577 vs 5510). Short-range negotiation finds
        fewer partners, but robots that hold for 10 steps widen their search (expanding
        ring), which breaks the deadlocks a short radius would otherwise cause.
  \item \textbf{Solvable instances on the main map}: 5--50 robots and up to 20\% obstacle
        density (1{,}700 runs of the main sweeps). No \texttt{local} run failed.
\end{itemize}

\subsection*{Summary}
The agents fail to complete their tasks when
(i) permanent disruptions cut goals off, even a single robot breaking down for good in a
one-cell aisle;
(ii) robots stuck in this way keep blocking others, which in these settings accounts for
more than half of the stuck robots;
(iii) a map is so crowded that one local hold cascades into a fleet-wide gridlock, the
case where global replanning does better;
(iv) obstacle density reaches about 30\%;
(v) robots cannot negotiate (\texttt{solo}).
Possible remedies are to detect unreachable goals and release or re-auction the task while
moving the robot aside, and to escalate to a larger regional replan once a hold cascade
exceeds a few robots.
